\begin{lexlean}{FoundryExperience}
\useglossary{lexlean.std.bool@1.2.0}
\useglossary{lexlean.std.nat@1.2.0}
\title{Boolean}
\begin{semanticmodule}
definition{
  name: oneLook,
  parameters: [],
  result: String,
  body: "Every surface Foundry draws itself, outside the renderer kit, takes the look of the rest of the app from the brand tokens in one stylesheet, in every theme: approval questions, notices of an open request, tool dialogs and the tools card. None writes a colour, a type face or a control style of its own. The result reads as clean, sharp and premium: one glass surface, a hairline edge, quiet text, no decoration."
}

definition{
  name: popupsCentred,
  parameters: [],
  result: String,
  body: "Every overlay popup is centred on the screen, across and down, over a light scrim. At most one approval question is shown at a time; the next waits its turn, and its time limit starts when it is shown."
}

definition{
  name: goldenRatio,
  parameters: [],
  result: String,
  body: "Every measure of these surfaces is from the golden ratio, 1.618: spaces from the series 5, 8, 13, 21, 34 and 55 pixels; a popup 550 pixels wide, which is 340 times the ratio, and at most 61.8 percent of the screen tall; 13 pixel text under a 16 pixel title; controls 34 pixels tall, and at least 44 under a finger."
}

definition{
  name: questionAsNotification,
  parameters: [],
  result: String,
  body: "A build-time question is drawn as the app draws its own notifications: a mark in a soft tile, a title, the explanation, the exact details behind one row, and two answers at the right. The refusal is plain and holds the focus; the approval is solid and is the last control."
}

definition{
  name: whyHowWhat,
  parameters: [],
  result: String,
  body: "Every explanation is ordered why, then how, then what: why this is being asked, how it happens and where the data goes, then what exactly is being approved and what cannot be undone. It is concise and self-descriptive, in plain words without technical jargon, and it still states every fact the approval boundaries require to be disclosed."
}

definition{
  name: sidebarIsNavigation,
  parameters: [],
  result: String,
  body: "The sidebar holds only navigation and is 233 pixels wide, which is 144 times the ratio."
}

definition{
  name: toolsOnSettings,
  parameters: [],
  result: String,
  body: "The tools are the first card of the Settings page: a title, then one quiet control each, with its own mark."
}

definition{
  name: homeAsksOneThing,
  parameters: [],
  result: String,
  body: "Home starts with the question of what to build, with no line under it and no app picker. The composer offers Build, Chat and Plan through its mode menu; the separate default-mode setting is not shown."
}

definition{
  name: homeIsAWorkspace,
  parameters: [],
  result: String,
  body: "Home is a workspace, not a page. Its line, Build your ideas with Prism, is one row at every size, the sentence in the page's own face and the name in a serif italic, so the line reads as a sentence and the name as a signature. The hero fills the first screen at every size, and the floor of apps already made peeks above the bottom edge, 89 pixels on a desktop and 55 on a phone, so one short scroll brings the cards and the way into the gallery into view, clear of the phone's tab bar. The hero sits at the golden section of the room above the floor on a phone as on a desktop, its question over a room at most 987 pixels wide; the box it is asked in is that room's golden section, 610 pixels, centred. At 768 pixels and above, two fresh ideas and Surprise me sit centred under it on one line; below 768 pixels, only Surprise me is shown. Surprise me writes another idea into the box. On a tall window the room under the hero is capped at 233 pixels, so it settles toward the middle instead of leaving a void. The floor is a panel that rises from the bottom edge, inset from the sides, on the popover glass under a hairline, with each app's name under its cover. On a phone the panel keeps the page gutter. The floor is looked at four ways, on one track at its left: Search, Featured, My apps and Bookmarked, the chosen one lit, with the way to all of them at its right. Featured is the same four Prism apps for every visitor, whether or not they have made anything; My apps previews what this browser holds; Bookmarked is what the visitor starred, kept in this browser alone and never sent anywhere. A card is its cover, a word on the cover saying whether it is published, then its name and who made it under its name. Cards share the panel's width in four even columns at 1024 pixels and above, and two below 1024 pixels, with no fixed card width. Search and Bookmarked select from the full local app collection before the four-card preview limit. In a mixed Bookmarked preview, featured apps keep their catalogue order before own apps in showcase order."
}

definition{
  name: composerRow,
  parameters: [],
  result: String,
  body: "The box's row reads left to right as the question is made: attach, then the model that will answer beside it, then the three ways to ask, then send. The way of asking is chosen on the same pill as the model, with its name and a chevron, and its menu lists each way on two lines, its name over a few words on what choosing it means. The names are the ones the kernel's catalogue gives: Build, Chat and Plan. One accessible trigger is normally reachable with Tab, names the current mode, and opens a keyboard-operated menu; there is no duplicate accessibility control."
}

definition{
  name: featuredCovers,
  parameters: [],
  result: String,
  body: "Featured apps use locally drawn spectrum artwork, white light entering and leaving as a spectrum turned to each app's own hue. An own app uses its checked, source-bound thumbnail when available, filling the cover from its top edge; otherwise it uses locally drawn spectrum artwork."
}

definition{
  name: modelMenu,
  parameters: [],
  result: String,
  body: "The model menu lists each model on two lines, its name over a few words on what choosing it means, with a tag saying where it runs, Cloud or Local, each with its mark, in one column at the right."
}
\end{semanticmodule}
\end{lexlean}
